% =====================================================================
%  Lecture 13 -- S-GW and P-GW                             (Persian edition)
% =====================================================================
\documentclass[11pt,a4paper]{article}
\providecommand\ACNroot{../..}
\input{\ACNroot/templates/lecture-fa}

\ACNset{number}{13}
\ACNset{title}{دروازه‌های خدمت‌دهنده و \lr{PDN} (\lr{S-GW} و \lr{P-GW})}
\ACNset{shorttitle}{دروازه‌های \lr{S-GW} و \lr{P-GW}}
\ACNset{titleother}{Serving and PDN Gateways (S-GW and P-GW)}
\ACNset{part}{بخش دوم -- دسترسی و هستهٔ نسل چهارم (درس‌های ۵ تا ۱۳)}
\ACNset{time}{۶ تا ۸ ساعت (۳ ساعت مطالعه، ۳ تا ۵ ساعت تمرین)}
\ACNset{level}{متوسط}
\ACNset{prereq}{درس‌های ۹ و ۱۲؛ نشانی‌دهی \lr{IP}، \lr{UDP}، تونل‌زنی}
\ACNset{gen}{4}

\begin{document}
\ACNcover

% ---------------------------------------------------------------------
\section{اهداف یادگیری}
پس از مطالعهٔ این درس باید بتوانید:
\begin{objectives}
  \item کارکردهای دروازهٔ خدمت‌دهنده و دروازهٔ \acro{PDN} را \textbf{از هم تمیز دهید} و \textbf{توضیح دهید} چرا صفحهٔ کاربر
        دو لنگر دارد.
  \item تونل‌زنی \lr{GTP-U} را \textbf{توصیف کنید} و سربار آن و اثرش بر \acro{MTU} را \textbf{محاسبه کنید}.
  \item تحویل فروسو به یک \acro{UE} بیکار و تغییر مسیر پس از تحویل را \textbf{دنبال کنید}.
  \item پلیس‌گذاری \lr{APN-AMBR} را با سطل نشانه \textbf{مدل کنید} و بافرها و نرخ بستهٔ دروازه را \textbf{اندازه‌گذاری کنید}.
  \item جداسازی صفحهٔ کنترل و کاربر (\acro{CUPS}) با قواعد \acro{PFCP} را \textbf{توضیح دهید} و رومینگ با مسیریابی خانگی و
        خروج محلی را \textbf{مقایسه کنید}.
\end{objectives}

% ---------------------------------------------------------------------
\section{دانش پیش‌نیاز}
این درس از مدل حامل درس ۹ (حامل‌های پیش‌فرض و اختصاصی، \acro{TFT}، \acro{QCI}، \acro{AMBR}) و رویه‌های \acro{MME} درس ۱۲
(اتصال، فراخوانی، درخواست خدمت) استفاده می‌کند. از شبکه به نشانی‌دهی \lr{IPv4/IPv6}، \acro{UDP} و ایدهٔ تونل‌زنی نیاز دارید:
حمل یک بسته به‌عنوان بار مفید بستهٔ دیگر.

\paragraph{واژه‌های کلیدی} واژگان این درس همراه با معادل انگلیسی آن‌ها:
\begin{keyterms}
دروازهٔ خدمت‌دهنده & \lr{S-GW} & لنگر صفحهٔ کاربر برای تحرک میان \acro{eNB}ها؛ بافر برای \acro{UE}های بیکار. \\
دروازهٔ \acro{PDN} & \lr{P-GW} & لنگر به سوی یک شبکهٔ دادهٔ بسته‌ای؛ نشانی \acro{IP} \acro{UE} را تخصیص می‌دهد. \\
پروتکل تونل‌زنی \acro{GPRS}، صفحهٔ کاربر & \lr{GTP-U} & پروتکل تونل‌زنی مبتنی بر \acro{UDP} برای بسته‌های کاربر \cite{ts29281}. \\
شناسهٔ نقطهٔ پایانی تونل & \lr{TEID} & شناسهٔ ۳۲ بیتی یک تونل در گرهٔ گیرنده. \\
تابع اعمال سیاست و صورت‌حساب & \lr{PCEF} & نقطهٔ اعمال قواعد \acro{PCC} در \lr{P-GW}. \\
اعلان دادهٔ فروسو & \lr{DDN} & درخواست \lr{S-GW} از \acro{MME} برای فراخوانی یک \acro{UE} بیکار. \\
نشانگر پایان & \lr{End marker} & بستهٔ \lr{GTP-U} که آخرین بسته روی مسیر قدیم پس از تغییر مسیر را نشان می‌دهد. \\
سطل نشانه & \lr{Token bucket} & مدل پلیس‌گذاری با نرخ $r$ و عمق $b$. \\
جداسازی صفحهٔ کنترل و کاربر & \lr{CUPS} & تقسیم هر دروازه به یک تابع کنترلی و یک تابع صفحهٔ کاربر \cite{ts23214}. \\
پروتکل کنترل پیش‌برد بسته & \lr{PFCP} & پروتکلی که صفحهٔ کنترل با آن قواعد بسته را نصب می‌کند \cite{ts29244}. \\
\end{keyterms}

% ---------------------------------------------------------------------
\section{مقدمه و انگیزه}
هر بایتی که یک مشترک \acro{LTE} می‌فرستد یا دریافت می‌کند از دو دروازه می‌گذرد. \emph{دروازهٔ خدمت‌دهنده} (\lr{S-GW}) رو به
شبکهٔ رادیویی است: در حالی که \acro{UE} میان \acro{eNB}ها جابه‌جا می‌شود ثابت می‌ماند و وقتی \acro{UE} خواب است بسته‌ها را
نگه می‌دارد. \emph{دروازهٔ \acro{PDN}} (\lr{P-GW}) رو به اینترنت یا \acro{IMS} است: به \acro{UE} نشانی \acro{IP} می‌دهد و
سیاست و صورت‌حساب اپراتور را اعمال می‌کند. \lr{P-GW} نقطهٔ حضور \acro{UE} در جهان \acro{IP} است، پس در طول یک نشست نباید
تغییر کند؛ \lr{S-GW} یک لنگر منطقه‌ای است و ممکن است با تحرک تغییر کند \cite{ts23401}.

دروازه‌ها جایی‌اند که هسته با حجم ترافیک روبه‌رو می‌شود. یک \acro{MME} بر حسب رویه در ثانیه اندازه‌گذاری می‌شود؛ یک دروازه
بر حسب بیت و بسته در ثانیه، و طراحی آن زیر سلطهٔ کارایی هدایت، بافرکردن و پردازش قواعد برای هر بسته است. این درس یک
بسته را در دروازه‌ها دنبال می‌کند، هزینهٔ تونل‌زنی و پلیس‌گذاری را کمّی می‌کند و با جداسازی صفحه‌های کنترل و کاربر که به
\acro{UPF} نسل پنجم می‌رسد به پایان می‌رسد.

% ---------------------------------------------------------------------
\section{مفاهیم بنیادی}

\subsection{تقسیم کار}
\begin{table}[htbp]
\centering\small
\caption{کارکردهای \lr{S-GW} و \lr{P-GW} \cite{ts23401}.}\label{tab:roles}
\rowcolors{2}{acnmist}{white}
\begin{tabularx}{\linewidth}{@{}Y Y@{}}
\toprule
\rowcolor{white}\lr{S-GW} & \lr{P-GW} \\
\midrule
لنگر محلی تحرک برای تحویل میان \acro{eNB}ها؛ ارسال نشانگرهای پایان & تخصیص نشانی \acro{IP} به \acro{UE} (\acro{IPv4}، پیشوند \acro{IPv6}) \\
بافرکردن فروسو در \lr{ECM-IDLE}؛ تحریک فراخوانی با \acro{DDN} & اعمال سیاست (\acro{PCEF}): دروازه‌بانی، کیفیت خدمات، \lr{APN-AMBR} \\
لنگر تحرک به سوی نسل ۲ و ۳ (\lr{S4}، \lr{S12}) & فیلترکردن بسته‌های فروسو با \acro{TFT}ها برای انتخاب حامل \\
شنود قانونی & شنود قانونی، بازرسی عمیق بسته \\
حساب‌داری هر حامل (مثلاً برای تسویهٔ رومینگ) & صورت‌حساب برخط (\lr{Gy}) و برون‌خط (\lr{Gz}) \\
علامت‌گذاری \acro{DSCP} روی شبکهٔ انتقال & علامت‌گذاری \acro{DSCP}، \acro{NAT} و دیوارهٔ آتش (به انتخاب اپراتور) \\
\bottomrule
\end{tabularx}
\end{table}

یک \acro{UE} در هر لحظه دقیقاً یک \lr{S-GW} دارد که به همهٔ اتصال‌های \acro{PDN} آن خدمت می‌دهد؛ ممکن است چند \lr{P-GW}
داشته باشد، یکی برای هر اتصال \acro{PDN} (مثلاً یکی برای \acro{APN} اینترنت و یکی برای \acro{IMS}). اپراتورها اغلب برای
ترافیک غیررومینگ یک \emph{\lr{S/P-GW} ترکیبی} مستقر می‌کنند که در آن \lr{S5} داخلی می‌شود.

\subsection{تونل به‌جای مسیر}
نشانی \acro{IP} \acro{UE} به مخزن نشانی \lr{P-GW} تعلق دارد. اگر مسیریاب‌های میان \acro{eNB} و \lr{P-GW} بسته‌های \acro{UE} را
بر اساس سرآیند \acro{IP} آن هدایت می‌کردند، هر تحویل به به‌روزرسانی مسیریابی نیاز داشت. در عوض هر حامل در تونل‌های
\lr{GTP-U} حمل می‌شود: \lr{S1-U} میان \acro{eNB} و \lr{S-GW} و \lr{S5/S8} میان \lr{S-GW} و \lr{P-GW} (\cref{fig:gw}). یک تحویل
فقط نقطهٔ پایانی تونل (نشانی و \acro{TEID}) را در \lr{S-GW} تغییر می‌دهد؛ نشانی \acro{UE} و \lr{P-GW} هرگز تغییر نمی‌کنند.
هر حامل جفت تونل خود را در هر جهت دارد که با \acro{TEID} گیرنده مشخص می‌شود.

\begin{acnfigure}
  \resizebox{\linewidth}{!}{\acnfig{lecture-13/gw-functions}}
  \caption{\lr{S-GW} و \lr{P-GW} روی مسیر صفحهٔ کاربر با کارکردهای اصلی آن‌ها.}
  \label{fig:gw}
\end{acnfigure}

% ---------------------------------------------------------------------
\section{توضیحات فنی تفصیلی}

\subsection{\lr{GTP-U}}
\lr{GTP-U} روی درگاه \acro{UDP} ۲۱۵۲ اجرا می‌شود \cite{ts29281}. سرآیند الزامی آن ۸ بایت است: پرچم‌ها (نسخهٔ ۱، نوع پروتکل
و بیت‌های \lr{E}، \lr{S} و \lr{PN} که ۴ بایت اختیاری برای سرآیند افزونه، شمارهٔ ترتیب و شمارهٔ \lr{N-PDU} را اعلام می‌کنند)،
نوع پیام (۲۵۵ برای \lr{G-PDU}، یعنی دادهٔ کاربر)، طول آنچه پس از سرآیند الزامی می‌آید و \acro{TEID} (\cref{fig:gtpu}).
پیام‌های سیگنالینگ خود \lr{GTP-U} شامل \lr{Echo Request/Response} برای نظارت مسیر، \lr{Error Indication} (گیرنده برای یک
\acro{TEID} زمینه‌ای ندارد) و \lr{End Marker} (نوع ۲۵۴) است. بستهٔ درونی (\lr{T-PDU}) تغییر نمی‌کند: دروازه‌ها سرآیند
\acro{IP} خود \acro{UE} را می‌بینند و پلیس‌گذاری می‌کنند.

\begin{acnfigure}
  \resizebox{\linewidth}{!}{\acnfig{lecture-13/gtpu-encap}}
  \caption{کپسوله‌سازی \lr{GTP-U} یک بستهٔ کاربر و سرآیند الزامی \lr{GTP-U}.}
  \label{fig:gtpu}
\end{acnfigure}

\subsection{تخصیص نشانی \lr{IP}}
هنگام اتصال (یا برقراری اتصال \acro{PDN}) \acro{UE} نوع \acro{PDN} را \acro{IPv4}، \acro{IPv6} یا \lr{IPv4v6} درخواست می‌کند.
برای \acro{IPv4}، \lr{P-GW} یک نشانی را در \lr{Create Session Response} برمی‌گرداند که در \acro{NAS} به \acro{UE} می‌رسد؛ برای
\acro{IPv6} یک پیشوند \lr{/64} تخصیص می‌دهد و آن را در \lr{Router Advertisement}ها روی حامل پیش‌فرض اعلام می‌کند و \acro{UE}
نشانی خود را از آن می‌سازد. نشانی‌های خصوصی \acro{IPv4} با \acro{NAT} در مقیاس اپراتور در \lr{P-GW} یا پشت آن رایج‌اند؛
این یکی از انگیزه‌های \acro{APN}های فقط \acro{IPv6} است.

\subsection{اعمال سیاست}
\acro{PCEF} در \lr{P-GW} قواعد \acro{PCC} دریافتی از \lr{Gx} (درس ۱۰) را اعمال می‌کند: بسته‌های فروسو را با \lr{DL TFT}ها به
حامل‌ها نگاشت می‌کند، جریان‌های دادهٔ خدمت را مسدود یا مجاز می‌کند (دروازه‌بانی)، \lr{GBR/MBR} هر حامل و
\lr{APN-AMBR} را روی حامل‌های غیر\acro{GBR} یک اتصال \acro{PDN} پلیس‌گذاری می‌کند (\cref{sec:tb})، حجم و زمان را برای هر
گروه نرخ‌گذاری جهت صورت‌حساب می‌شمارد و بسته‌ها را با مقادیر \acro{DSCP} برگرفته از \acro{QCI} علامت‌گذاری می‌کند تا شبکهٔ
انتقال بتواند آن‌ها را اولویت‌بندی کند. \lr{UE-AMBR} را \acro{eNB} اعمال می‌کند.

\subsection{حالت بیکار و تحویل}
وقتی \acro{UE} به \lr{ECM-IDLE} می‌رود، \acro{eNB} زمینهٔ آن را آزاد می‌کند و \acro{MME} از \lr{S-GW} می‌خواهد تونل‌های
\lr{S1-U} را آزاد کند؛ بخش \lr{S5/S8} باقی می‌ماند. آن‌گاه یک بستهٔ فروسو هیچ تونل \lr{S1-U} نمی‌یابد: \lr{S-GW} آن را بافر
می‌کند و \acro{DDN} را به \acro{MME} می‌فرستد که \acro{UE} را فراخوانی می‌کند (\cref{fig:ddn}). پس از یک تحویل \lr{X2}،
\acro{eNB} مقصد از \acro{MME} تغییر مسیر می‌خواهد؛ \acro{MME} پیام \lr{Modify Bearer Request} را به \lr{S-GW} می‌فرستد که
فروسو را به \acro{eNB} مقصد می‌برد و یک یا چند بستهٔ \emph{نشانگر پایان} روی مسیر قدیم می‌فرستد تا \acro{eNB} مبدأ بداند
هدایت بسته‌های در راه کی تمام شده است و مقصد بتواند بسته‌ها را به ترتیب تحویل دهد.

\subsection{جداسازی صفحهٔ کنترل و کاربر (\lr{CUPS})}
انتشار ۱۴ هر دروازه را به یک تابع کنترلی (\lr{SGW-C}، \lr{PGW-C}) و یک یا چند تابع صفحهٔ کاربر (\lr{SGW-U}، \lr{PGW-U})
تقسیم کرد که با واسط‌های \lr{Sx} به هم وصل‌اند \cite{ts23214}. تابع کنترلی \lr{S11}، \lr{S5-C}، \lr{Gx} و \lr{Gy} را نگه
می‌دارد؛ تابع صفحهٔ کاربر فقط بسته‌ها را بر اساس قواعدی که با \acro{PFCP} نصب می‌شوند هدایت می‌کند \cite{ts29244}:
\begin{itemize}
  \item \textbf{\acro{PDR}} (قاعدهٔ تشخیص بسته): قاعده بر کدام بسته‌ها (\acro{TEID}، نشانی \acro{UE}، فیلتر \acro{SDF}) اعمال می‌شود؛
  \item \textbf{\acro{FAR}} (قاعدهٔ اقدام هدایت): هدایت، دورریختن، بافرکردن یا تکثیر، و به کجا؛
  \item \textbf{\acro{QER}} (قاعدهٔ اعمال کیفیت خدمات): دروازه‌بانی و محدودیت‌های نرخ؛
  \item \textbf{\acro{URR}} (قاعدهٔ گزارش مصرف): چه چیزی شمرده و کی گزارش شود؛
  \item \textbf{\acro{BAR}} (قاعدهٔ اقدام بافر): نحوهٔ بافرکردن برای \acro{UE}های بیکار.
\end{itemize}
آن‌گاه توابع صفحهٔ کاربر را می‌توان نزدیک شبکهٔ رادیویی قرار داد (برای تأخیر کم و ترافیک محلی) و مستقل از صفحهٔ کنترل
مقیاس داد. \acro{UPF} نسل پنجم و واسط \lr{N4} آن همین مدل را به کار می‌برند (درس ۲۳).

% ---------------------------------------------------------------------
\section{مدل‌های ریاضی}

\subsection{سربار تونل‌زنی و \lr{MTU}}\label{sec:gtp}
\paragraph{مدل} یک بستهٔ کاربر با $L$ بایت با سربار تونل $H$ (\acro{IP} بیرونی + \acro{UDP} + \lr{GTP-U}) حمل می‌شود:
$H = 20+8+8 = 36$ بایت با انتقال \acro{IPv4} و $H = 56$ بایت با \acro{IPv6}. بازدهی انتقال و بزرگ‌ترین بستهٔ کاربر که بدون
قطعه‌قطعه‌شدن در \acro{MTU} انتقال $M$ جا می‌شود عبارت‌اند از
\begin{equation}\label{eq:gtpeff}
  \eta = \frac{L}{L+H}, \qquad L_{\max} = M - H .
\end{equation}
\paragraph{مثال} $L = 1400$ با انتقال \acro{IPv4}: $\eta = 97.5\,\%$. با $M = 1500$، $L_{\max} = 1464$ بایت (۱۴۴۴ با انتقال
\acro{IPv6})؛ به همین دلیل اپراتورها \acro{MTU} \acro{UE} را کمتر از ۱۵۰۰ پیکربندی می‌کنند یا \acro{MTU} انتقال را بزرگ‌تر
می‌کنند.
\paragraph{تفسیر} سربار برای بسته‌های کوچک مهم است: برای صدا کسر بزرگی از بسته است (\cref{ex:volte}). فشرده‌سازی سرآیند
(\acro{ROHC}) روی رادیو در هسته کمکی نمی‌کند، چون آنجا سرآیندهای کامل و تونل حمل می‌شوند.

\subsection{بافرکردن فروسو}\label{sec:buffer}
\paragraph{مدل} یک منبع با نرخ $R$ به یک \acro{UE} بیکار ارسال می‌کند؛ تا تونل \lr{S1-U} پس از تأخیر فراخوانی و درخواست
خدمت $D$ بازیابی شود، \lr{S-GW} باید این مقدار را بافر کند:
\begin{equation}\label{eq:buffer}
  B = R\,D .
\end{equation}
\paragraph{مثال} $R = 2$ مگابیت بر ثانیه و $D = 1.5$ ثانیه: $B = 3$ مگابیت $= 375$ کیلوبایت. در عمل منابع \acro{TCP} پیش از
انتظار برای تصدیق فقط چند بسته می‌فرستند، پس بافرهای واقعی بسیار کوچک‌ترند و دروازه‌ها بافر هر حامل را محدود و مازاد را
دور می‌ریزند.

\subsection{پلیس‌گذاری با سطل نشانه}\label{sec:tb}
\paragraph{مدل} یک سطل نشانه با نرخ $r$ و عمق $b$ یک جریان ترافیک را می‌پذیرد اگر برای هر بازه به طول $T$، دادهٔ ارسالی
$A(T)$ در شرط زیر صدق کند:
\begin{equation}\label{eq:tb}
  A(T) \le b + r\,T .
\end{equation}
منبعی با نرخ اوج $R > r$ که با سطل پر آغاز می‌کند می‌تواند به مدت $T^\ast = b/(R-r)$ با نرخ اوج ارسال کند.
\paragraph{مثال} \lr{APN-AMBR} $r = 10$ مگابیت بر ثانیه و $b = 8$ مگابیت (۱ مگابایت): در ۳ ثانیه حداکثر $8 + 30 = 38$
مگابیت (۴٫۷۵ مگابایت) عبور می‌کند؛ یک منبع ۴۰ مگابیتی می‌تواند به مدت $8/30 = 0.27$ ثانیه انفجاری ارسال کند و سپس به ۱۰
مگابیت بر ثانیه محدود می‌شود (\cref{fig:tb}).
\paragraph{تفسیر} نرخ سهم بلندمدت را تعیین می‌کند؛ عمق تعیین می‌کند ترافیک چقدر می‌تواند انفجاری باشد. سطل خیلی کوچک به
\acro{TCP} آسیب می‌زند (اتلاف در هر انفجار) و سطل خیلی بزرگ اجازه می‌دهد انفجارهای کوتاه به‌طور محسوسی از نرخ مشترک‌شده
فراتر روند.

\begin{acnfigure}
  \acnfig{lecture-13/token-bucket}
  \caption{پوش سطل نشانه $b + rT$ (\cref{eq:tb}) و یک منبع با نرخ اوج.}
  \label{fig:tb}
\end{acnfigure}

\subsection{نرخ بسته}\label{sec:pps}
\paragraph{مدل} هزینهٔ هدایت به ازای هر بسته است، نه هر بایت. دروازه‌ای که گذردهی $C$ با اندازهٔ میانگین بستهٔ $\bar L$ را حمل
می‌کند، این تعداد بسته در ثانیه را پردازش می‌کند:
\begin{equation}\label{eq:pps}
  P = \frac{C}{8\,\bar L}\ \text{packets/s}.
\end{equation}
\paragraph{مثال} $C = 40$ گیگابیت بر ثانیه و $\bar L = 800$ بایت: $P = 6.25$ میلیون بسته در ثانیه -- هر بسته باید در برابر
\acro{PDR}های نشست خود طبقه‌بندی، شمارش و کپسوله شود، به‌طور میانگین در حدود ۱۶۰ نانوثانیه برای هر بسته در یک خط
پردازش واحد؛ به همین دلیل صفحه‌های کاربر از ورودی/خروجی بسته با دورزدن هسته و تخلیه به سخت‌افزار استفاده می‌کنند.

% ---------------------------------------------------------------------
\section{معماری شبکه و نمودارها}
در یک شبکهٔ غیررومینگ، \lr{S-GW} و \lr{P-GW} معمولاً در مراکز دادهٔ منطقه‌ای کنار هم‌اند؛ با \acro{CUPS}، توابع صفحهٔ کاربر
به سوی لبه حرکت می‌کنند و توابع کنترلی مرکزی می‌مانند. در رومینگ، شبکهٔ میزبان \lr{S-GW} را فراهم می‌کند؛ \lr{P-GW} یا در
شبکهٔ خانگی است (\emph{مسیریابی خانگی}، از طریق \lr{S8}، حالت معمول برای داده) یا در شبکهٔ میزبان (\emph{خروج محلی}، مثلاً
برای صدای \acro{IMS} در برخی توافق‌ها). مسیریابی خانگی سیاست و صورت‌حساب را در شبکهٔ خانگی نگه می‌دارد اما مسیر را
طولانی‌تر می‌کند (\cref{exr:roaming}).

% ---------------------------------------------------------------------
\section{رویه‌های گام‌به‌گام}

\begin{procedurebox}{دادهٔ فروسو برای یک \lr{UE} بیکار}\label{proc:ddn}
\begin{steps}
  \item یک بستهٔ فروسو به \lr{P-GW} می‌رسد که آن را با \lr{DL TFT} به یک حامل نگاشت و از طریق \lr{S5} به \lr{S-GW}
        هدایت می‌کند.
  \item \lr{S-GW} برای حامل تونل \lr{S1-U} ندارد: بسته را بافر و \lr{Downlink Data Notification} را به \acro{MME} می‌فرستد
        (بسته‌های بعدی بدون \acro{DDN} جدید بافر می‌شوند).
  \item \acro{MME} \acro{UE} را در فهرست \acro{TAI} آن فراخوانی می‌کند؛ \acro{UE} با \lr{Service Request} پاسخ می‌دهد.
  \item \acro{MME} زمینه را در \acro{eNB} برقرار می‌کند (\lr{Initial Context Setup})، \acro{TEID} \acro{eNB} را می‌آموزد و
        \lr{Modify Bearer Request} را به \lr{S-GW} می‌فرستد.
  \item \lr{S-GW} بسته‌های بافرشده و سپس ترافیک جدید را به \acro{eNB} هدایت می‌کند.
\end{steps}
\end{procedurebox}

\begin{acnfigure}
  \resizebox{\linewidth}{!}{\acnfig{lecture-13/ddn}}
  \caption{درخواست خدمت آغازشده از شبکه با بافرکردن در \lr{S-GW} (\cref{proc:ddn}).}
  \label{fig:ddn}
\end{acnfigure}

\begin{procedurebox}{تغییر مسیر پس از تحویل \lr{X2}}\label{proc:switch}
\begin{steps}
  \item در طول تحویل، \acro{eNB} مبدأ بسته‌های فروسوی در راه را از طریق \lr{X2} به \acro{eNB} مقصد هدایت می‌کند.
  \item \acro{eNB} مقصد \lr{Path Switch Request} را با \acro{TEID}های فروسوی خود به \acro{MME} می‌فرستد.
  \item \acro{MME} \lr{Modify Bearer Request} را به \lr{S-GW} می‌فرستد که تونل‌های فروسو را به مقصد می‌برد.
  \item \lr{S-GW} بسته‌های نشانگر پایان را روی تونل‌های قدیم به \acro{eNB} مبدأ می‌فرستد؛ مبدأ آن‌ها را به مقصد رله می‌کند
        و مقصد بسته‌های هدایت‌شده را پیش از بسته‌های جدید تحویل می‌دهد.
  \item \acro{eNB} مبدأ زمینهٔ \acro{UE} را آزاد می‌کند.
\end{steps}
\end{procedurebox}

% ---------------------------------------------------------------------
\section{مثال‌های حل‌شده}

\begin{examplebox}{سربار یک بستهٔ صوتی}\label{ex:volte}
\textbf{مسئله.} یک فریم \lr{AMR-WB} با نرخ ۱۲٫۶۵ کیلوبیت بر ثانیه هر ۲۰ میلی‌ثانیه حدود ۳۳ بایت بار مفید \acro{RTP} تولید
می‌کند؛ \acro{UE} از \acro{IPv6} استفاده می‌کند. اندازهٔ بسته روی \lr{S1-U} (انتقال \acro{IPv4}) و کسری را که سرآیند است
محاسبه کنید.

\textbf{حل.} بستهٔ درونی: $33 + 12\ (\mathrm{RTP}) + 8\ (\mathrm{UDP}) + 40\ (\mathrm{IPv6}) = 93$ بایت. روی \lr{S1-U}:
$93 + 36 = 129$ بایت. سربار تونل به‌تنهایی $36/129 = 28\,\%$ است؛ همهٔ سرآیندها با هم $(129-33)/129 = 74\,\%$. با ۵۰ بسته
در ثانیه، تماس روی \lr{S1-U} برای ۱۲٫۶۵ کیلوبیت بر ثانیه گفتار $129\times8\times50 = 51.6$ کیلوبیت بر ثانیه مصرف می‌کند.
\end{examplebox}

\begin{examplebox}{پذیرش با سطل نشانه}\label{ex:tb}
\textbf{مسئله.} با $r = 10$ مگابیت بر ثانیه و $b = 8$ مگابیت، کاربری یک پروندهٔ ۵ مگابایتی را با ۴۰ مگابیت بر ثانیه از یک
کارساز دریافت می‌کند و با سطل پر آغاز می‌کند. دریافت چقدر طول می‌کشد؟

\textbf{حل.} $5$ مگابایت $= 40$ مگابیت. در انفجار اولیهٔ $T^\ast = 0.27$ ثانیه، $40\times0.27 = 10.7$ مگابیت عبور می‌کند؛
$29.3$ مگابیت باقی‌مانده با ۱۰ مگابیت بر ثانیه در $2.93$ ثانیه عبور می‌کند. مجموع $\approx 3.2$ ثانیه، در برابر ۱ ثانیه بدون
پلیس‌گذاری و ۴ ثانیه با $b = 0$.
\end{examplebox}

\begin{examplebox}{نرخ بستهٔ دروازه}\label{ex:pps}
\textbf{مسئله.} یک \lr{P-GW} به $2\times10^5$ کاربر هم‌زمان فعال با میانگین ۲۰۰ کیلوبیت بر ثانیه خدمت می‌دهد. گذردهی و نرخ
بسته را برای بستهٔ میانگین ۸۰۰ بایتی محاسبه کنید.

\textbf{حل.} $C = 2\times10^5\times2\times10^5 = 40$ گیگابیت بر ثانیه؛ $P = 40\times10^9/6400 = 6.25$ میلیون بسته در ثانیه
(\cref{eq:pps}). بسته‌های فراسو (بسیاری تصدیق‌های کوچک \acro{TCP}) به نرخ بسته می‌افزایند در حالی که گذردهی اندکی اضافه
می‌کنند.
\end{examplebox}

% ---------------------------------------------------------------------
\section{مقایسهٔ فناوری‌ها}
\begin{table}[htbp]
\centering\small
\caption{دروازه‌های بسته‌ای در نسل‌های مختلف.}\label{tab:cmp}
\rowcolors{2}{acnmist}{white}
\begin{tabularx}{\linewidth}{@{}l Y Y Y@{}}
\toprule
\rowcolor{white} & \lr{SGSN + GGSN} (نسل ۲ و ۳) & \lr{S-GW + P-GW} (نسل ۴) & \lr{UPF (+ SMF)} (نسل ۵) \\
\midrule
لنگرها            & \acro{GGSN} (\acro{IP})، \acro{SGSN} (تحرک) & \lr{P-GW} (\acro{IP})، \lr{S-GW} (تحرک) & \acro{UPF} لنگر نشست \acro{PDU}، \acro{UPF}های میانی \\
تونل‌زنی          & \lr{GTP-U} (\lr{Gn}، \lr{Iu-PS})  & \lr{GTP-U} (\lr{S1-U}، \lr{S5/S8}) & \lr{GTP-U} (\lr{N3}، \lr{N9}) \\
کنترل هدایت       & یکپارچه                          & یکپارچه؛ \acro{CUPS} از انتشار ۱۴ & همیشه جدا: \acro{PFCP} روی \lr{N4} \\
دانه‌بندی کیفیت خدمات & زمینهٔ \acro{PDP}            & حامل \acro{EPS}                   & جریان کیفیت خدمات (\acro{QFI}) \\
خروج محلی / لبه   & نادر                             & با \acro{CUPS} ممکن               & بومی (طبقه‌بند فراسو، نقطهٔ انشعاب) \\
\bottomrule
\end{tabularx}
\end{table}

\paragraph{یک لنگر یا دو لنگر؟} دو لنگر دغدغه‌ها را جدا می‌کنند: \lr{S-GW} تحرک رادیویی را از \lr{P-GW} پنهان می‌کند و
\lr{P-GW} شبکهٔ \acro{PDN} را از سمت رادیو. هزینهٔ آن یک گام و حالت اضافی است. نسل پنجم این ایده را نگه می‌دارد اما آن را
انعطاف‌پذیر می‌کند: یک نشست ممکن است یک \acro{UPF} یا زنجیره‌ای از آن‌ها داشته باشد که \acro{SMF} برای هر نشست انتخاب
می‌کند.

% ---------------------------------------------------------------------
\section{نکته‌های کلیدی}
\begin{takeawaybox}
\begin{itemize}
  \item \lr{S-GW}: لنگر تحرک میان \acro{eNB}ها، بافرکردن و \acro{DDN} برای \acro{UE}های بیکار، نشانگرهای پایان.
  \item \lr{P-GW}: لنگر \acro{IP} (تخصیص نشانی)، \acro{PCEF} (نگاشت \acro{TFT}، دروازه‌بانی، \acro{AMBR}، صورت‌حساب)، \lr{SGi}
        به \acro{PDN}.
  \item تونل‌های \lr{GTP-U} حامل‌ها را حمل می‌کنند؛ سربار $H = 36$ (\acro{IPv4}) یا ۵۶ (\acro{IPv6}) بایت، $L_{\max} = M - H$.
  \item بافرکردن $B = RD$؛ پلیس‌گذاری $A(T) \le b + rT$؛ هزینهٔ هدایت $P = C/(8\bar L)$ بسته در ثانیه.
  \item \acro{CUPS} دروازه‌ها را به صفحه‌های کنترل و کاربر تقسیم می‌کند که با قواعد \acro{PFCP} (\lr{PDR, FAR, QER, URR, BAR})
        به هم وصل‌اند -- نیای مستقیم \acro{UPF} نسل پنجم.
\end{itemize}
\end{takeawaybox}

% ---------------------------------------------------------------------
\section{تمرین‌ها}

\begin{exercisebox}{1}{کدام دروازه؟}\label{exr:roles}
در هر حالت کدام دروازه (\lr{S-GW} یا \lr{P-GW}) تغییر می‌کند: (الف) تحویل \lr{X2} میان دو \acro{eNB} در ناحیهٔ یک \lr{S-GW}؛
(ب) \acro{UE} به منطقه‌ای می‌رود که \lr{S-GW} دیگری به آن خدمت می‌دهد؛ (ج) \acro{UE} اتصال \acro{PDN} دومی به \acro{APN}
\acro{IMS} باز می‌کند؟
\end{exercisebox}

\begin{exercisebox}{1}{نشانگرهای پایان}\label{exr:endmarker}
توضیح دهید اگر \lr{S-GW} بسته‌های نشانگر پایان را نمی‌فرستاد، پس از تحویل در \acro{eNB} مقصد چه مشکلی ممکن بود پیش آید.
\end{exercisebox}

\begin{exercisebox}{2}{\lr{MTU}}\label{exr:mtu}
شبکهٔ انتقال \acro{MTU} برابر ۱۵۰۰ بایت دارد و از \acro{IPv6} استفاده می‌کند. بزرگ‌ترین بستهٔ \acro{UE} که از قطعه‌قطعه‌شدن
پرهیز می‌کند چیست؟ شبکهٔ انتقال برای حمل بسته‌های ۱۵۰۰ بایتی \acro{UE} به چه \acro{MTU}ای نیاز دارد؟
\end{exercisebox}

\begin{exercisebox}{2}{اندازه‌گذاری بافر}\label{exr:buffer}
یک \lr{S-GW} بافر هر حامل را به ۲۵۶ کیلوبایت محدود می‌کند. برای تأخیر فراخوانی و درخواست خدمت ۱٫۲ ثانیه، تا چه نرخ منبعی
هیچ بسته‌ای دور ریخته نمی‌شود؟ اگر $10^4$ \acro{UE} هم‌زمان فراخوانی شوند و هر بافر پر شود، چقدر حافظه لازم است؟
\end{exercisebox}

\begin{exercisebox}{3}{رومینگ با مسیریابی خانگی}\label{exr:roaming}
مشترکی ۴۰۰۰ کیلومتر دور از شبکهٔ خانگی رومینگ می‌کند و به کارسازی در کشور میزبان دسترسی می‌یابد. تأخیر رفت‌وبرگشت
اضافی مسیریابی خانگی را در مقایسه با خروج محلی تخمین بزنید (انتشار با $2\times10^8$ متر بر ثانیه، بدون پردازش) و بررسی
کنید چرا اپراتورها هنوز برای داده مسیریابی خانگی را ترجیح می‌دهند.
\end{exercisebox}

\begin{exercisebox}{3}{مقیاس‌دهی صفحهٔ کاربر}\label{exr:cups}
یک هستهٔ \lr{PGW-U} ۱٫۵ میلیون بسته در ثانیه را هدایت می‌کند. اگر هسته‌ها حداکثر تا ۷۰٪ بارگذاری شوند، برای ۱۰۰ گیگابیت
بر ثانیه با بستهٔ میانگین ۷۰۰ بایتی چند هسته لازم است؟ وقتی ترافیک دو برابر می‌شود اما تعداد مشترکین نه، \acro{CUPS} به
اپراتور اجازهٔ چه کاری را می‌دهد؟
\end{exercisebox}

\subsection*{پاسخ‌ها}
\addcontentsline{toc}{subsection}{پاسخ‌ها}

\begin{solutionbox}{\ref{exr:roles}}
(الف) هیچ‌کدام: فقط نقطهٔ پایانی \lr{S1-U} (\acro{eNB}) تغییر می‌کند. (ب) \lr{S-GW} تغییر می‌کند (جابه‌جایی \lr{S-GW})؛
\lr{P-GW} و نشانی \acro{IP} می‌مانند. (ج) یک \lr{P-GW} جدید (احتمالاً) برای اتصال \acro{PDN} جدید افزوده می‌شود؛ \lr{S-GW}
مشترک است.
\end{solutionbox}

\begin{solutionbox}{\ref{exr:endmarker}}
مقصد بسته‌ها را از دو مسیر دریافت می‌کند: هدایت‌شده از طریق \lr{X2} از مبدأ و مستقیم از \lr{S-GW}. بدون نشانگر پایان
نمی‌تواند بداند جریان هدایت‌شده کی تمام شده است؛ تحویل نخست بسته‌های جدید جریان را بی‌ترتیب می‌کند (به \acro{TCP} آسیب
می‌زند) و انتظار با زمان‌سنج تأخیر می‌افزاید. نشانگر پایان مسیر قدیم را به‌صراحت می‌بندد.
\end{solutionbox}

\begin{solutionbox}{\ref{exr:mtu}}
$L_{\max} = 1500 - 56 = 1444$ بایت. برای حمل بسته‌های ۱۵۰۰ بایتی \acro{UE}، \acro{MTU} انتقال باید دست‌کم
$1500 + 56 = 1556$ بایت باشد (یا ۱۵۴۰ با انتقال \acro{IPv4})؛ اپراتورها اغلب در شبکهٔ بک‌هال از قاب‌های بزرگ استفاده
می‌کنند.
\end{solutionbox}

\begin{solutionbox}{\ref{exr:buffer}}
$R \le B/D = 256\times8\ \text{kbit}/1.2\ \text{s} = 1.71$ مگابیت بر ثانیه. حافظهٔ بدترین حالت $10^4\times256$ کیلوبایت
$= 2.56$ گیگابایت؛ در عمل بسیار کمتر، چون منابع \acro{TCP} پس از یک پنجره متوقف می‌شوند و بیشتر انفجارها چند بسته‌اند.
\end{solutionbox}

\begin{solutionbox}{\ref{exr:roaming}}
مسیریابی خانگی هر بسته را میزبان $\to$ خانه $\to$ کارساز (در کشور میزبان) و به همان راه برمی‌گرداند: $2\times4000$ کیلومتر
اضافی در هر جهت، یعنی $4\times4000\ \text{km}/(2\times10^8\ \text{m/s}) = 80$ میلی‌ثانیه رفت‌وبرگشت اضافی. اپراتورها آن را
می‌پذیرند چون شبکهٔ خانگی کنترل سیاست، صورت‌حساب، شنود قانونی و دسترسی به خدمات (مثلاً \acro{APN}های خانگی) را حفظ
می‌کند و تسویهٔ رومینگ ساده‌تر است؛ خروج محلی به توافق‌هایی دربارهٔ سیاست و صورت‌حساب در شبکهٔ میزبان نیاز دارد.
\end{solutionbox}

\begin{solutionbox}{\ref{exr:cups}}
$P = 100\times10^9/(8\times700) = 17.9$ میلیون بسته در ثانیه؛ $17.9/1.5 = 11.9$ هسته در بار کامل و $11.9/0.7 = 17.0 \to 18$
هسته. با \acro{CUPS}، اپراتور ظرفیت صفحهٔ کاربر (نمونه‌ها یا هسته‌های بیشتر \lr{PGW-U}، احتمالاً در سایت‌های لبهٔ جدید) را
بدون دست‌زدن به صفحهٔ کنترل می‌افزاید، که بار آن به مشترکین و رویه‌ها بستگی دارد نه به حجم ترافیک.
\end{solutionbox}

% ---------------------------------------------------------------------
\section{مطالعهٔ بیشتر}
کارکردهای دروازه‌ها در \lr{TS 23.401} \cite{ts23401}، \lr{GTPv2-C} در \lr{TS 29.274} \cite{ts29274} و \lr{GTP-U} در
\lr{TS 29.281} \cite{ts29281} مشخص شده‌اند؛ \acro{CUPS} و \acro{PFCP} در \lr{TS 23.214} \cite{ts23214} و \lr{TS 29.244}
\cite{ts29244}؛ کنترل سیاست و صورت‌حساب در \lr{TS 23.203} \cite{ts23203}. السون و همکاران \cite{olsson2013epc} استقرار
دروازه‌ها و رومر و همکاران \cite{rommer2020core} تکامل به سوی صفحهٔ کاربر نسل پنجم را بررسی می‌کنند.

\begin{latin}
\printbibliography[heading=subbibliography,title={References}]
\end{latin}
\end{document}
